\documentclass[tikz,border=3mm,multi=tikzpicture]{standalone}
\usepackage{eleve-quimica}

\begin{document}

% FIGURA 1 — fig-01-equilibrio-dinamico-graficos
\begin{tikzpicture}[quimica figura, x=1cm, y=1cm]
  \path[use as bounding box] (-0.15,-0.15) rectangle (10.65,12.0);

  % Velocidades.
  \draw[quimica eixo] (1.25,6.75) -- (10.25,6.75)
    node[above left=2pt, quimica rotulo] {tempo};
  \draw[quimica eixo] (1.25,6.75) -- (1.25,11.35)
    node[above=2pt, quimica rotulo] {velocidade};
  \draw[quimica auxiliar] (6.65,6.75) -- (6.65,11.1);
  \node[quimica rotulo, below=6pt] at (6.65,6.75) {$t_{eq}$};

  \draw[quimica direta]
    (1.55,10.75) .. controls (3.0,9.25) and (4.8,8.75) .. (6.65,8.75)
    -- (9.75,8.75);
  \draw[quimica inversa]
    (1.55,7.15) .. controls (3.0,8.35) and (4.8,8.75) .. (6.65,8.75)
    -- (9.75,8.75);
  \node[quimica rotulo, text=quimicaazul, anchor=west] at (2.10,10.30)
    {$v_{\mathrm{direta}}$};
  \node[quimica rotulo, text=quimicalaranja, anchor=west] at (2.10,7.35)
    {$v_{\mathrm{inversa}}$};
  \node[quimica destaque, anchor=east] at (9.75,9.18)
    {$v_{\mathrm{direta}}=v_{\mathrm{inversa}}$};

  % Concentrações.
  \draw[quimica eixo] (1.25,0.85) -- (10.25,0.85)
    node[above left=2pt, quimica rotulo] {tempo};
  \draw[quimica eixo] (1.25,0.85) -- (1.25,5.45)
    node[above=2pt, quimica rotulo] {concentração};
  \draw[quimica auxiliar] (6.65,0.85) -- (6.65,5.2);
  \node[quimica rotulo, below=6pt] at (6.65,0.85) {$t_{eq}$};

  \draw[quimica direta]
    (1.55,4.95) .. controls (3.0,3.75) and (4.8,3.20) .. (6.65,3.20)
    -- (9.75,3.20);
  \draw[quimica inversa]
    (1.55,1.20) .. controls (3.0,2.85) and (4.8,4.10) .. (6.65,4.10)
    -- (9.75,4.10);
  \node[quimica rotulo, text=quimicaazul, anchor=east] at (9.75,2.78)
    {$[\mathrm{N_2O_4}]$};
  \node[quimica rotulo, text=quimicalaranja, anchor=east] at (9.75,4.55)
    {$[\mathrm{NO_2}]$};
\end{tikzpicture}


% FIGURA 2 — fig-02-pressao-na-sintese-da-amonia
\begin{tikzpicture}[quimica figura, x=1cm, y=1cm]
  \path[use as bounding box] (-0.15,-0.15) rectangle (10.8,10.2);

  \node[quimica rotulo] at (5.15,9.75) {reagentes: 4 mols gasosos};
  \draw[quimica recipiente] (0.75,6.00) rectangle (9.55,9.25);

  \node[quimica processo, minimum width=1.45cm] at (1.85,7.62)
    {$\mathrm{N_2}$};
  \node[quimica processo, minimum width=1.45cm] at (4.05,7.62)
    {$\mathrm{H_2}$};
  \node[quimica processo, minimum width=1.45cm] at (6.25,7.62)
    {$\mathrm{H_2}$};
  \node[quimica processo, minimum width=1.45cm] at (8.45,7.62)
    {$\mathrm{H_2}$};

  \draw[quimica fluxo] (5.15,5.72) -- (5.15,4.68);
  \node[quimica destaque, anchor=west] at (5.55,5.18)
    {aumento da pressão};

  \node[quimica rotulo] at (5.15,4.05)
    {produtos: 2 mols gasosos — lado favorecido};
  \draw[quimica recipiente] (0.75,0.55) rectangle (9.55,3.55);
  \node[quimica separacao, minimum width=2.10cm] at (3.45,2.05)
    {$\mathrm{NH_3}$};
  \node[quimica separacao, minimum width=2.10cm] at (6.85,2.05)
    {$\mathrm{NH_3}$};
\end{tikzpicture}


% FIGURA 3 — fig-03-fluxo-haber-bosch
\begin{tikzpicture}[quimica figura, x=1cm, y=1cm]
  \path[use as bounding box] (-0.20,-0.25) rectangle (11.7,9.7);

  \node[quimica rotulo] (feed) at (0.95,7.25)
    {$\mathrm{N_2+3H_2}$};
  \node[quimica processo, minimum width=2.05cm] (compressor) at (3.15,7.25)
    {compressão};
  \node[quimica processo, minimum width=2.65cm] (reator) at (6.30,7.25)
    {reator\\Fe; $T$ moderada};
  \node[quimica separacao, minimum width=2.55cm] (separador) at (9.65,7.25)
    {resfriamento\\e separação};

  \draw[quimica fluxo] (feed.east) -- (compressor.west);
  \draw[quimica fluxo] (compressor.east) -- (reator.west);
  \draw[quimica fluxo] (reator.east) -- (separador.west);

  \node[quimica separacao, minimum width=2.25cm] (produto) at (9.65,3.90)
    {$\mathrm{NH_3}$ retirada};
  \draw[quimica fluxo, draw=quimicaverde] (separador.south) -- (produto.north);

  \coordinate (retorno1) at (9.65,1.55);
  \coordinate (retorno2) at (3.15,1.55);
  \draw[quimica retorno] (separador.east) -- ++(0.85,0)
    |- (retorno1) -- (retorno2) -- (compressor.south);
  \node[quimica destaque] at (6.30,1.12)
    {$\mathrm{N_2}$ e $\mathrm{H_2}$ não convertidos retornam};
\end{tikzpicture}

\end{document}
